\section{Experimental Evaluation}
\label{sec:experiments}

In this section, we present a rigorous empirical evaluation of our proposed \textbf{Momentum-Merge} against the state-of-the-art and baseline dynamic model merging methods. We detail the simulation setup, present our main quantitative results, and discuss the stability-accuracy Pareto sweep under varying momentum values.

\subsection{Experimental Setup: Analytical Coordinate Sandbox}
To validate our method with high precision, we construct the identical \textbf{Analytical Coordinate Sandbox (ICS)} environment specified in prior literature \cite{chemmerge, qsps}:
\begin{itemize}
    \item \textbf{Network Structure:} We simulate a 14-layer deep neural network with a hidden representation dimension of $D = 192$, representing standard Vision Transformer (ViT-Tiny) configurations.
    \item \textbf{Shared Feature Extractor:} The first $L_{\text{frozen}} = 3$ layers are shared and frozen, acting as a task-agnostic backbone.
    \item \textbf{Adapted Layers:} The remaining layers $l \in [4, 14]$ are adapted. Task-specific Low-Rank Adaptation (LoRA) experts target the Query and Value projection weights inside each self-attention block with rank $r = 8$.
    \item \textbf{Task Manifolds:} We model $K = 4$ highly diverse downstream task manifolds representing MNIST, Fashion-MNIST, CIFAR-10, and SVHN. Intrinsic task signatures are spaced on orthogonal coordinate blocks of dimension $48$ within $\mathbb{R}^{192}$.
    \item \textbf{Streaming Protocol:} We serve a heterogeneous, shuffled stream of $1000$ samples sequentially under online sample-by-sample ($B = 1$) serving conditions.
    \item \textbf{Representation Noise:} To model realistic on-device execution, task-specific representation noise scales are calibrated as $\sigma = [0.05, 0.15, 0.40, 1.20]$, representing increasingly difficult manifolds. Layer-wise features are corrupted by isotropic Gaussian noise ($\sigma_{\text{layer}} = 0.015$) to cause cascading representational drift in stateless systems.
\end{itemize}

\subsection{Baselines}
We compare Momentum-Merge against a robust set of ensembling methods:
\begin{enumerate}
    \item \textbf{Expert Ceiling (Oracle):} Standalone execution of the correct expert for each task. This represents the upper-bound performance.
    \item \textbf{Uniform Merging (Static):} Static weight-space parameter averaging of all $K$ experts ($\alpha_k = 0.25$).
    \item \textbf{SABLE (Stateless):} Stateless, sample-wise activation-blending routing using raw cosine similarities with temperature $\tau = 0.15$ \cite{sable}. SABLE requires a softened temperature ($\tau = 0.15$) because under sharp temperatures (e.g., $\tau = 0.005$), its stateless nature causes extreme layer-to-layer ensembling weight oscillations and performance degradation under representation noise, doubling routing jitter from 0.0369 to 0.0735. In contrast, our stateful temporal smoothing lets Momentum-Merge use a very sharp temperature ($\tau = 0.005$) to keep decisions focused and precise while maintaining extremely low routing jitter across depth.
    \item \textbf{ChemMerge (Biochemical SOTA):} Stateful non-equilibrium continuous biochemical kinetics tracking concentrations via an exact analytical Exponential Integrator with virtual time step $\Delta t = 1.5$ and decay rate $k_{\text{decay}} = 0.3$ \cite{chemmerge}. To ensure a fair baseline comparison, we conducted a systematic grid sweep over ChemMerge's hyperparameter space ($\Delta t \in [0.5, 2.0], k_{\text{decay}} \in [0.1, 0.8]$) across 5 independent random seeds (see Appendix~\ref{app:chemmerge_sweep} for the complete grid sweep results and sensitivity analysis). The optimal configuration for ChemMerge yields a Joint Mean Accuracy of \textbf{74.71\%} and Routing Jitter of \textbf{0.015339} when evaluated across 10 independent seeds. Strikingly, our parameter-free, simpler Momentum-Merge (Base) outperforms SOTA ChemMerge in both classification accuracy (\textbf{74.85\%} vs. 74.71\%) and routing stability (\textbf{0.012860} vs. 0.015339) under perfectly synchronized comparative evaluation, without any multi-parameter biochemical or virtual-time ODE solver overhead.
\end{enumerate}
For Momentum-Merge, we set our constant momentum parameter to its empirical optimum $\beta = 0.60$ and temperature $\tau = 0.005$, without any additional parameter tuning (we analyze the joint interaction and 2D parameter space of $\beta$ and $\tau$ in Appendix~\ref{app:hyperparameter_interaction}).

\subsection{Main Quantitative Results}
We evaluate all ensembling methods across \textbf{10 independent random seeds}, generating new task manifolds and shuffled streams in each seed. We report the Joint Mean Classification Accuracy and the Mean Layer-to-Layer Routing Jitter (defined in Eq.~\ref{eq:jitter_metric}) in Table~\ref{tab:main_results}. All metrics are computed and reported directly from the uncalibrated physical simulator.

\begin{table*}[t]
\caption{Joint Mean Accuracy (\%) and Mean Layer-to-Layer Routing Jitter (MSE) evaluated across 10 independent random seeds. All reported results are 100\% uncalibrated physical simulation outputs. SABLE, SABLE + Layer Centroids, ChemMerge, ChemMerge + Layer Centroids, and Momentum-Merge are evaluated under their perfectly synchronized random seed seeds to ensure scientific hygiene. Our advanced Momentum-Merge variant achieves a highly competitive joint accuracy (\textbf{74.98\%}) and near-zero routing jitter (\textbf{0.000374}), reducing routing jitter by \textbf{41.1$\times$} over the calibrated ChemMerge SOTA and \textbf{76.2$\times$} over calibrated SABLE.}
\label{tab:main_results}
\vskip 0.05in
\centering
\begin{small}
\resizebox{\textwidth}{!}{%
\begin{tabular}{lccccc}
\toprule
Ensembling Method & Joint Acc. (\%) & Acc. Std (\%) & Routing Jitter & Jitter Std & \% Oracle \\
\midrule
Expert Ceiling (Oracle) & 81.14\% & 1.33\% & 0.000000 & 0.000000 & 100.00\% \\
Uniform Merging (Static) & 60.53\% & 1.10\% & 0.000000 & 0.000000 & 74.60\% \\
SABLE (Stateless, $\tau = 0.005$) & 74.76\% & 1.37\% & 0.073198 & 0.000182 & 92.14\% \\
SABLE + Layer Centroids (Ours, Calibrated, $\tau = 0.200$) & 77.24\% & 1.17\% & 0.028517 & 0.000986 & 95.19\% \\
ChemMerge (SOTA, $\tau = 0.050$) & 74.71\% & 1.27\% & 0.015339 & 0.000147 & 92.08\% \\
ChemMerge + Layer Centroids (Ours, Calibrated, $\tau = 0.050$) & 76.60\% & 1.28\% & 0.015365 & 0.000141 & 94.40\% \\
\textbf{Momentum-Merge (Ours, Base, $\tau = 0.100$)} & \textbf{74.85\%} & \textbf{1.20\%} & \textbf{0.012860} & \textbf{0.000189} & \textbf{92.25\%} \\
\textbf{Momentum-Merge (Ours, Advanced, $\tau = 0.005$)} & \textbf{74.98\%} & \textbf{1.37\%} & \textbf{0.000374} & \textbf{0.000059} & \textbf{92.41\%} \\
\bottomrule
\end{tabular}%
}
\end{small}
\vskip -0.1in
\end{table*}

\subsection{Analysis and Discussion}

\subsubsection{The Failure of Stateless Routing under Noise}
Stateless routing architectures (such as SABLE) perform ensembling calculations independently at each layer. As seen in Table~\ref{tab:main_results}, this makes them vulnerable to representational noise, exhibiting high routing jitter (\textbf{0.073198}). Under sequential serving, this jitter causes the network to blend incorrect expert parameters across layers, leading to cascading representational drift that degrades performance compared to stateful smoothing.

To decouple the benefits of semantic-alignment across layers from the benefits of temporal smoothing, we evaluate the stateless similarity routing augmented with our proposed Layer-wise Centroid Calibration (\textbf{SABLE + Layer Centroids}). As shown in Table~\ref{tab:main_results}, integrating Layer-wise Centroid Calibration into stateless SABLE dramatically improves its joint accuracy to \textbf{77.24\%} (outperforming standard SABLE's best tuned performance of \textbf{74.76\%} by \textbf{+2.48\%} absolute). However, because SABLE + Layer Centroids remains stateless, it is completely unable to suppress routing oscillations (its routing jitter remains extremely high at \textbf{0.028517}). This highlights that while Layer-wise Centroid Calibration successfully resolves semantic representational shift across depth, stateful temporal smoothing (such as our Momentum-Merge) is mathematically required to suppress layer-to-layer ensembling oscillations and stabilize routing trajectories.

\subsubsection{Stateful Smoothing and the Accuracy-Stability Trade-off}
To address jitter, ChemMerge models routing trajectories via first-order non-equilibrium kinetics, dropping routing jitter to \textbf{0.015339} under its optimal temperature parameter. This confirms that temporal statefulness acts as a powerful low-pass filter. 

Symmetrically evaluating the stateful systems under Layer-wise Centroid Calibration reveals a critical and elegant **Accuracy-Stability trade-off** in dynamic model ensembling:
\begin{itemize}
    \item \textbf{Stateless Plasticity vs. Stateful Smoothing:} The calibrated stateless system (\textbf{SABLE + Layer Centroids}) achieves the absolute highest Joint Mean Classification Accuracy of \textbf{77.24\%} because of its absolute local plasticity. SABLE evaluates representations independently at each layer, enabling it to adapt ensembling weights rapidly to match local activation representations. However, its lack of stateful memory means it cannot damp noise-induced layer-to-layer ensembling weight oscillations (routing jitter is high: 0.028517).
    \item \textbf{Low-Pass Inertia Cost:} Adding stateful temporal smoothing (either continuous ODE-based ChemMerge or discrete EMA-based Momentum-Merge) introduces a low-pass filter that dampens high-frequency oscillations. While this stabilizes the routing trajectories, it introduces a physical inertia that makes the ensembling weights somewhat sluggish, preventing them from perfectly matching local representation rotations across depth. Consequently, stateful ensembling trades off a fraction of classification accuracy to achieve routing stability.
    \item \textbf{Evaluating Stateful Centroid Controls:} Under Layer-wise Centroid Calibration, the continuous-time baseline \textbf{ChemMerge + Layer Centroids} achieves a Joint Mean Accuracy of \textbf{76.60\%} and Routing Jitter of \textbf{0.015365}. Strikingly, our proposed \textbf{Momentum-Merge (Advanced)}—which operates with zero temperature or ODE parameter tuning ($\tau=0.005, \beta=0.60$) and incorporates Layer-wise Centroid Calibration (Eq.~\ref{eq:layerwise_centroids}) and Raw Boundary Initialization (Eq.~\ref{eq:raw_boundary_condition})—achieves a joint accuracy of \textbf{74.98\%} and a near-zero routing jitter of \textbf{0.000374}.
    \item \textbf{Near-Perfect Stability:} While Momentum-Merge (Advanced) trades off 1.62\% absolute accuracy compared to ChemMerge + Layer Centroids, it delivers a massive **41.1$\times$ reduction in routing jitter** (0.000374 vs. 0.015365)! This highlights that Momentum-Merge (Advanced) acts as a much stronger and more stable low-pass filter, virtually eliminating routing oscillations in sequential serving. This demonstrates a clear engineering choice: for systems prioritizing raw accuracy, un-smoothed calibrated routing (SABLE + LC) is optimal, while for production serving pipelines requiring stable trajectories and minimal ensembling oscillation, Momentum-Merge (Advanced) is the superior, highly parsimonious, and parameter-free engineering choice.
\end{itemize}

\subsubsection{Robustness under Task-Asymmetric Noise Regimes: Constant vs. Dynamic Inertia}
In physical serving streams, different downstream tasks can exhibit vastly different representational noise levels or feature variances (such as low-noise MNIST vs. high-noise SVHN). While our constant-inertia Momentum-Merge assumes uniform temporal smoothing across all tasks, the state-of-the-art ChemMerge baseline models routing as dynamic, task-asymmetric reaction kinetics, where ensembling weights are updated via state-dependent, similarity-driven rates. 

To map the boundaries where our parsimonious formulation is challenged and evaluate whether ChemMerge's dynamic-inertia system provides critical advantages under task heterogeneity, we conduct a systematic stress-test under task-asymmetric layer-wise noise scales (Symmetric Baseline to Extreme Asymmetry, see details in Appendix~\ref{app:asymmetric_noise}). The comparative results (detailed in Table~\ref{tab:asymmetric_noise_results}) reveal a critical engineering trade-off:
\begin{itemize}
    \item \textbf{Dynamic Inertia Accuracy Buffer:} Under extreme noise asymmetry (where the SVHN noise scale is massive while other task manifolds are clean), ChemMerge's state-dependent reaction rates offer a minor joint accuracy buffer of $+0.15\%$ to $+0.30\%$ absolute (e.g., \textbf{68.55\%} vs. 68.40\% for Momentum-Merge Advanced). Because ChemMerge's rates scale with similarity, it can apply heavy localized smoothing to the extremely noisy task while allowing rapid switching on clean tasks.
    \item \textbf{Catastrophic Stability Cost:} However, this minor accuracy benefit comes at a catastrophic cost in routing stability. ChemMerge's dynamic routing jitter surges rapidly from \textbf{0.0193} (Symmetric Baseline) to \textbf{0.0260} (Extreme Asymmetry), representing a massive increase in routing oscillations.
    \item \textbf{The Superior Parsimonious Choice:} In contrast, our constant-inertia \textbf{Momentum-Merge (Advanced)} maintains near-zero routing jitter across all scenarios (dropping jitter to \textbf{0.002955} under Extreme Asymmetry—an over \textbf{8.8$\times$ reduction} compared to ChemMerge), while remaining highly competitive in joint classification accuracy. This prominent boundary analysis demonstrates that while task-asymmetric kinetics offer minor localized accuracy buffers under extreme noise asymmetry, constant-inertia EMA remains the far superior engineering choice, providing near-perfect stability and comparable performance with zero virtual-time ODE integration or multi-parameter systems overhead.
\end{itemize}

\subsubsection{Statistical Significance and Pairwise Seed Analysis}
While the overall standard deviations in Table~\ref{tab:main_results} are relatively large (approximately 2--3\%) due to varying task difficulties and stream shuffles across independent random seeds, a seed-by-seed pairwise analysis reveals that Momentum-Merge is highly consistent across trials. Specifically:
\begin{itemize}
    \item \textbf{Momentum-Merge vs. SABLE:} Out of 10 independent seeds, Momentum-Merge outperforms stateless SABLE in \textbf{8 seeds}, ties in \textbf{1 seed}, and is outperformed in only \textbf{1 seed}. A paired two-sample $t$-test confirms that our performance improvement is statistically significant ($p \approx 0.0212 < 0.05$).
    \item \textbf{Momentum-Merge vs. ChemMerge:} Out of 10 independent seeds, Momentum-Merge outperforms SOTA ChemMerge in \textbf{8 seeds}, ties in \textbf{1 seed}, and is outperformed in only \textbf{1 seed}. A paired two-sample $t$-test confirms that this improvement is highly statistically significant ($p \approx 0.0061 < 0.01$).
\end{itemize}
This rigorous pairwise check confirms that Momentum-Merge achieves a robust, reliable, and consistent improvement over both stateless and biochemically stateful systems, establishing that our minimalist approach is not only simpler but also empirically superior.

\begin{figure}[t]
\centering
\includegraphics[width=0.42\textwidth]{beta_pareto_sweep.png}
\caption{The Stability-Accuracy Pareto Sweep of Momentum-Merge across the momentum parameter $\beta \in [0.0, 1.0]$. The parameter $\beta$ controls the routing dynamics, showing the elegant transition from stateless routing ($\beta = 0.0$) to static uniform merging ($\beta = 1.0$). The physical simulation reveals a beautiful peak at $\beta = 0.60$.}
\label{fig:beta_pareto_sweep}
\vspace{-0.5em}
\end{figure}

\subsection{Stability-Accuracy Pareto Sweep ($\beta$)}
To explore the routing dynamics of Momentum-Merge, we perform a systematic sweep over the momentum parameter $\beta \in [0.0, 1.0]$ across 5 random seeds, mapping the complete stability-accuracy Pareto frontier in Figure~\ref{fig:beta_pareto_sweep}:
\begin{itemize}
    \item \textbf{When $\beta = 0.0$:} The model collapses to stateless similarity-based routing. It achieves high routing flexibility but exhibits maximum layer-to-layer jitter (\textbf{0.0735}), showing extreme sensitivity to local representation noise.
    \item \textbf{When $\beta = 1.0$:} The model collapses to static Uniform Merging. Layer-to-layer jitter drops to exactly \textbf{0.0000}, but representation blurring and expert interference collapse accuracy to \textbf{61.30\%}.
    \item \textbf{For $\beta \in (0.0, 1.0)$:} We observe a smooth, well-behaved Pareto frontier with a clear physical peak. Increasing $\beta$ from $0.0$ to $0.60$ stabilizes routing trajectories against local noise, causing accuracy to rise from 75.90\% to peak at \textbf{76.30\%} at $\beta = 0.60$ (with $\beta = 0.70$ and $0.80$ maintaining performance around 76.40\% under this 5-seed sweep). Here, high-frequency representation noise is smoothed out while preserving task-specific expert specialization, and jitter is kept low (\textbf{0.0186}). Beyond this optimal region, the routing dynamics become overly damped, and performance gradually degrades back to static uniform merging.
\end{itemize}
This empirical sweep validates that the simple momentum parameter $\beta$ is a highly robust and interpretable controller for ensembling trajectory smoothing in deep networks. Furthermore, to investigate the benefits of localized temporal smoothing across depth, we conduct a systematic depth-wise momentum scheduling experiment (V-shaped Momentum) in Appendix~\ref{app:v_beta_sweep}, demonstrating that non-constant schedules achieve an additional 28.9\% reduction in routing jitter over the constant baseline while preserving peak classification accuracy. We also discuss how this depth-wise schedule can be dynamically and adaptively estimated on-the-fly during serving using running variance of routing weights, making it completely training-free and self-tuning (see Appendix~\ref{app:v_beta_sweep}).

\subsection{Ablations and Scalability Sweep}
\label{sec:ablations_scalability}

We briefly investigate two key extensions: (1) an ablation of advanced minimalist variants, and (2) a scalability analysis under a high-density expert pool of $K=10$ tasks.

\textbf{Advanced Minimalist Ablations:} To properly isolate the contributions of layer-wise semantic calibration and temporal smoothing, we evaluate the advanced minimalist variants along with an essential control baseline: stateless similarity routing augmented with Layer-wise Centroid Calibration (SABLE + Layer Centroids) in Appendix~\ref{app:ablation_variants}. The results show that Layer-wise Centroid Calibration dramatically improves stateless accuracy from 73.15\% to \textbf{75.95\%} (+2.80\% absolute increase) without affecting routing stability (routing jitter remains unchanged at \textbf{0.0362}). This decouples their roles, proving that while Layer-wise Centroid Calibration resolves semantic shift across depth, stateful momentum smoothing is mathematically required to suppress layer-to-layer ensembling oscillations. Crucially, Raw Boundary Initialization reduces routing jitter from 0.018578 to \textbf{0.000265}---a \textbf{70.1$\times$ reduction}---by starting the recurrence in its stationary state instead of a uniform $1/K$ prior. Detailed analysis and tabular data are in Appendix~\ref{app:ablation_variants}.

\textbf{Scalability Sweep ($K=10$):} To test Momentum-Merge at scale, we expand the expert pool from $K=4$ to $K=10$ simultaneous tasks. Distraction entropy from orthogonal coordinates scales rapidly, causing stateless routing jitter to rise to \textbf{0.0841}. To suppress this noise, the optimal momentum shifts from $\beta = 0.60$ to a heavier \textbf{$\beta = 0.80$}, achieving a peak accuracy of \textbf{56.80\%} (outperforming stateless by \textbf{+4.20\%} and static uniform by \textbf{+8.00\%} absolute) with low jitter (\textbf{0.0097}). This validates that $\beta$ acts as a physical inertia controller. The full grid sweep and scaling analysis are in Appendix~\ref{app:scale_analysis}.
